% ============================================================================
%  CAPÍTULO II: OBJETIVOS, ALCANCE Y LÍMITES
% ============================================================================

\section{Objetivos del trabajo}

\subsection{Objetivo general}

Implementar, sobre la instalaci\'on de \textbf{Debian 12 (Bookworm)} del
servidor \texttt{dc1}, el \textbf{DNS del dominio} \texttt{sudoers.lan}
administrado por el Controlador de Dominio de \textbf{Samba 4} (zona
autoritativa y registros de descubrimiento) y un \textbf{servidor web} basado
en \textbf{Nginx} que publique el portal del proyecto y act\'ue como
\textit{proxy inverso} de los servicios de la intranet, accesible por nombre
de dominio y con cifrado \textbf{HTTPS} desde cualquier estaci\'on de la red.

\subsection{Objetivos espec\'ificos}

\begin{enumerate}[label=\textbf{\arabic*.}]
  \item Confirmar que el DC de \textbf{Samba 4} sirve como \textbf{DNS
        autoritativo} del dominio: zona \texttt{sudoers.lan} con el backend
        \textbf{SAMBA\_INTERNAL}, registros A de los servicios y registros
        SRV de descubrimiento (LDAP, Kerberos).
  \item Crear los \textbf{registros A} de los servicios publicados
        (\texttt{web}, \texttt{www}, \texttt{print}, \texttt{portainer},
        \texttt{joaquin}, \texttt{david}, \texttt{nicolas}) apuntando al
        servidor \texttt{192.168.0.2}, de forma idempotente mediante el
        script \texttt{dc/dns-records.sh} y \texttt{samba-tool dns}.
  \item Configurar el \textbf{reenv\'io hacia Internet} (\textit{forwarder})
        hacia el router \texttt{192.168.0.1} para que el mismo DNS resuelva
        nombres externos.
  \item Desplegar \textbf{Nginx} como contenedor Docker con el portal
        est\'atico de la intranet y los \textit{virtual hosts} de
        \texttt{sudoers.lan}, \texttt{print.sudoers.lan} y
        \texttt{portainer.sudoers.lan}.
  \item Implementar \textbf{HTTPS} con \textbf{CA interna} y certificado
        \textit{wildcard} \texttt{*.sudoers.lan}, con redirecci\'on autom\'atica
        de HTTP (80) a HTTPS (443).
  \item Exponer por el \textit{proxy} los \textbf{subdominios de los equipos}
        de los administradores (\texttt{david.sudoers.lan},
        \texttt{nicolas.sudoers.lan}) hacia las m\'aquinas
        \texttt{192.168.0.51:8080} y \texttt{192.168.0.52:8080}.
  \item Validar el conjunto con pruebas de resoluci\'on (\texttt{dig},
        \texttt{samba-tool dns query}) y de acceso web (\texttt{curl},
        navegador), desde el propio servidor y desde un cliente de la red.
\end{enumerate}